\documentclass[aspectratio=169]{beamer}

% Theme selection
\usetheme{Madrid}
\usecolortheme{default}

% Packages
\usepackage[utf8]{inputenc}
\usepackage[T1]{fontenc}
\usepackage{graphicx}
\usepackage{booktabs}
\usepackage{amsmath}
\usepackage{amssymb}
\usepackage{hyperref}

% Custom colors (optional)
% \definecolor{myblue}{RGB}{0,82,155}
% \setbeamercolor{structure}{fg=myblue}

% Title page information
\title{Your Presentation Title}
\subtitle{Optional Subtitle}
\author{Your Name}
\institute{Your Institution \\ Your Department}
\date{\today}

% Footer customization
\setbeamertemplate{footline}{
  \leavevmode%
  \hbox{%
  \begin{beamercolorbox}[wd=.333333\paperwidth,ht=2.25ex,dp=1ex,center]{author in head/foot}%
    \usebeamerfont{author in head/foot}\insertshortauthor
  \end{beamercolorbox}%
  \begin{beamercolorbox}[wd=.333333\paperwidth,ht=2.25ex,dp=1ex,center]{title in head/foot}%
    \usebeamerfont{title in head/foot}\insertshorttitle
  \end{beamercolorbox}%
  \begin{beamercolorbox}[wd=.333333\paperwidth,ht=2.25ex,dp=1ex,right]{date in head/foot}%
    \usebeamerfont{date in head/foot}\insertshortdate{}\hspace*{2em}
    \insertframenumber{} / \inserttotalframenumber\hspace*{2ex} 
  \end{beamercolorbox}}%
  \vskip0pt%
}

\begin{document}

% Title slide
\begin{frame}
\titlepage
\end{frame}

% Outline
\begin{frame}{Outline}
\tableofcontents
\end{frame}

% Section 1
\section{Introduction}

\begin{frame}{Introduction}
\begin{itemize}
    \item Introduce your topic
    \item Provide context and background
    \item State the problem or research question
\end{itemize}
\end{frame}

\begin{frame}{Motivation}
\begin{block}{Why is this important?}
Explain the significance of your work and why the audience should care.
\end{block}

\begin{itemize}
    \item Key point 1
    \item Key point 2
    \item Key point 3
\end{itemize}
\end{frame}

\begin{frame}{Research Questions}
\begin{enumerate}
    \item First research question
    \item Second research question
    \item Third research question
\end{enumerate}
\end{frame}

% Section 2
\section{Background}

\begin{frame}{Background}
\begin{columns}
\column{0.5\textwidth}
\textbf{Key Concept 1}
\begin{itemize}
    \item Detail 1
    \item Detail 2
\end{itemize}

\column{0.5\textwidth}
\textbf{Key Concept 2}
\begin{itemize}
    \item Detail 1
    \item Detail 2
\end{itemize}
\end{columns}
\end{frame}

\begin{frame}{Related Work}
\begin{table}
\centering
\begin{tabular}{lcc}
\toprule
Approach & Metric 1 & Metric 2 \\
\midrule
Method A & 0.85 & 0.78 \\
Method B & 0.87 & 0.80 \\
\textbf{Our Method} & \textbf{0.92} & \textbf{0.85} \\
\bottomrule
\end{tabular}
\caption{Comparison with existing methods}
\end{table}
\end{frame}

% Section 3
\section{Methodology}

\begin{frame}{Proposed Approach}
\begin{itemize}
    \item Component 1: Description
    \item Component 2: Description
    \item Component 3: Description
\end{itemize}

\begin{center}
% Uncomment when you have a figure
% \includegraphics[width=0.6\textwidth]{figures/approach-diagram.pdf}
\fbox{\parbox{0.6\textwidth}{\centering [Your diagram will be placed here]}}
\end{center}
\end{frame}

\begin{frame}{Mathematical Formulation}
Consider the optimization problem:

\begin{equation}
\min_{x} f(x) \quad \text{subject to} \quad g(x) \leq 0
\end{equation}

Where:
\begin{itemize}
    \item $f(x)$ is the objective function
    \item $g(x)$ represents constraints
\end{itemize}
\end{frame}

\begin{frame}{Implementation}
\begin{block}{Key Implementation Details}
\begin{enumerate}
    \item Step 1: Data preprocessing
    \item Step 2: Model training
    \item Step 3: Evaluation
\end{enumerate}
\end{block}
\end{frame}

% Section 4
\section{Results}

\begin{frame}{Experimental Setup}
\begin{itemize}
    \item \textbf{Datasets}: Description of datasets used
    \item \textbf{Metrics}: Evaluation metrics
    \item \textbf{Baselines}: Comparison methods
\end{itemize}
\end{frame}

\begin{frame}{Main Results}
\begin{center}
% Uncomment when you have a figure
% \includegraphics[width=0.8\textwidth]{figures/results-chart.pdf}
\fbox{\parbox{0.8\textwidth}{\centering [Your results chart will be placed here]}}
\end{center}

\begin{itemize}
    \item Our method achieves X\% improvement
    \item Significant in Y metric
\end{itemize}
\end{frame}

\begin{frame}{Detailed Analysis}
\begin{columns}
\column{0.5\textwidth}
\textbf{Strengths}
\begin{itemize}
    \item[+] Advantage 1
    \item[+] Advantage 2
    \item[+] Advantage 3
\end{itemize}

\column{0.5\textwidth}
\textbf{Limitations}
\begin{itemize}
    \item[-] Limitation 1
    \item[-] Limitation 2
\end{itemize}
\end{columns}
\end{frame}

% Section 5
\section{Conclusion}

\begin{frame}{Summary}
\begin{block}{Key Takeaways}
\begin{enumerate}
    \item Main finding 1
    \item Main finding 2
    \item Main finding 3
\end{enumerate}
\end{block}
\end{frame}

\begin{frame}{Contributions}
\begin{itemize}
    \item \textbf{Contribution 1}: Brief description
    \item \textbf{Contribution 2}: Brief description
    \item \textbf{Contribution 3}: Brief description
\end{itemize}
\end{frame}

\begin{frame}{Future Work}
\begin{enumerate}
    \item Extend the approach to...
    \item Investigate the impact of...
    \item Apply the method to...
\end{enumerate}
\end{frame}

% Thank you slide
\begin{frame}[plain]
\centering
\Huge Thank You!

\vspace{1cm}

\Large Questions?

\vspace{1cm}

\normalsize
\texttt{your.email@institution.edu}
\end{frame}

% References slide (optional)
\begin{frame}[allowframebreaks]{References}
\tiny
\begin{thebibliography}{99}
\bibitem{ref1} Author, A. (2023). Title of paper. \textit{Journal Name}, 10(2), 123-145.
\bibitem{ref2} Author, B. (2022). Title of book. Publisher.
% Add more references as needed
\end{thebibliography}
\end{frame}

\end{document}
